\documentclass{article}
\usepackage{graphicx} % Required for inserting images
\usepackage{amsmath} % Essential for symbols like the double-struck R
\usepackage{amsfonts} % Needed for \mathbb

\title{BohdansLegacyLatex}

\begin{document}

\maketitle

\begin{description}

\item{(a)} \[
\text{Min} \left\{ J(x) = -\sum_{i=1}^{m} v^i \ln \left( g(x^T u^i) \right) + (1 - v^i) \ln \left( 1 - g(x^T u^i) \right)  \mid x \in \mathbb{R}^5 \right\}
\]
Dosadime za kazde $g(z)$ funkciu ktorej sa rovna $g(z) = \frac{1}{1+e^{-z}}$
\[ \ln \left( g(x^Tu^i) \right) =\ln \left( \frac{1}{1+e^{-x^Tu^i}} \right) =-\ln \left( 1+e^{-x^Tu^i} \right) \]

\[ \ln \left( 1-g(x^Tu^i) \right) =\ln \left( 1-\frac{1}{1+e^{-x^Tu^i}} \right) = \] \[= \ln \left( \frac{e^{-x^Tu^i}}{1+e^{-x^Tu^i}} \right) =\ln \left( e^{-x^Tu^i} \right) -\ln \left( 1+e^{-x^Tu^i} \right) \]

Teraz prepiseme to co mame vnutri sumy 
\[ v^i \left( -\ln \left( 1+e^{-x^Tu^i} \right) \right)+(1-v^i) \left( \ln \left( e^{-x^Tu^i} \right) -\ln \left( 1+e^{-x^Tu^i} \right) \right) = \] \[= (1-v^i)\ln \left( e^{-x^Tu^i} \right)-\ln \left( 1+e^{-x^Tu^i} \right)=(1-v^i)\left( -x^Tu^i \right)-\ln \left( 1+e^{-x^Tu^i} \right) \]
Zanesieme minus do vnutra sumy a dostaneme
\[ (1-v^i)x^Tu^i+\ln \left( 1+e^{-x^Tu^i} \right) \]
Vo vysledku dostaneme
\[
\text{Min} \left\{ J(x) = \sum_{i=1}^{m} (1-v^i)x^Tu^i+\ln \left( 1+e^{-x^Tu^i} \right)  \mid x \in \mathbb{R}^5 \right\},
\]

\item({b})
Zderivujeme najprv podla kazdej zlozky vektora x
\[ \frac{\delta J}{\delta x_j}=\sum_{i=1}^{m}(x^Tu^i-v^ix^Tu^i+\ln(1+e^{-x^Tu^i}))' = \] \[= \sum_{i=1}^{m}(u_j^i - v^iu_j^i+\frac{e^{-x^Tu^i}}{1+e^{-x^Tu^i}}(-u_j^i))=\sum_{i=1}^{m}((1-v^i-\frac{e^{-x^Tu^i}}{1+e^{-x^Tu^i}})u_j^i) = \] \[= \sum_{i=1}^{m}((\frac{1+e^{-x^Tu^i}-e^{-x^Tu^i}}{1+e^{-x^Tu^i}}-v^i)u_j^i)=\sum_{i=1}^{m}((\frac{1}{1+e^{-x^Tu^i}}-v^i)u_j^i) \]
Teraz mozeme vyjadrit gradient
\[ \nabla J(x) = \sum_{i=1}^{m}((\frac{1}{1+e^{-x^Tu^i}}-v^i)u^i)\]

\end{description}

\end{document}
